\section{Cadre méthodologique}
\label{sec:methodology}

\subsection{Moteur de backtest}

L'ensemble des calculs repose sur \textbf{VectorBT Pro} (v2026.3.1), qui offre trois avantages clés pour ce type de recherche~:
\begin{itemize}
\item \textbf{Pipelines Numba-compilés} pour les grid searches (résolution minute sur $\approx$ 3~millions de barres par pair).
\item \textbf{Portefeuilles natifs multi-symbole} avec \code{cash\_sharing} et groupement, permettant de partager un pool de trésorerie entre les paires d'un même sleeve.
\item \textbf{Chunking mémoire} pour les sweeps de paramètres, avec appels explicites au \code{vbt.flush()} entre chunks~--- critique pour tenir dans l'enveloppe RAM lors des benchmarks récents (Phase~17, optimisation \emph{perf/chunking-ram}).
\end{itemize}

\subsection{Données}

\begin{center}
\renewcommand{\arraystretch}{1.15}
\begin{tabular}{lll}
\toprule
\textbf{Source} & \textbf{Fréquence} & \textbf{Usage} \\
\midrule
EUR-USD, GBP-USD, USD-JPY, USD-CAD & minute & MR Macro (close+vwap) \\
Idem (resamplés) & daily & TS Momentum, RSI Daily \\
Spread 10Y--2Y US Treasury & daily & Filtre macro MR Macro \\
Unemployment rate US & mensuelle & Filtre macro MR Macro \\
\bottomrule
\end{tabular}
\end{center}

Les parquets FX sont strictement intersectés sur leur index minute (pas de \code{ffill} entre paires) et le pipeline rejette toute intersection qui rétrécit de plus de 5\,\% par rapport à la plus grande série disponible.

\subsection{Découpage in-sample / out-of-sample}

\begin{itemize}
\item \textbf{In-sample}~: 1\textsuperscript{er} avril 2019 $\rightarrow$ 1\textsuperscript{er} avril 2025, soit 2260 barres daily (~6 ans).
\item \textbf{Out-of-sample}~: 1\textsuperscript{er} avril 2025 $\rightarrow$ 1\textsuperscript{er} avril 2026, soit 314 barres (~14\,\% de l'IS).
\end{itemize}

\begin{warningbox}
\textbf{Point critique de transparence}~: la fenêtre OOS a été choisie \emph{ex-post}~--- l'utilisateur a demandé la vérification out-of-sample en Phase~18, lorsque les données 2025-04+ étaient déjà disponibles. Ce n'est donc pas un vrai \emph{holdout} pré-engagé. Malgré ce caveat, le résultat est suffisamment fort pour être significatif (Sharpe OOS 1.44 $>$ Sharpe IS 0.94).
\end{warningbox}

\subsection{Anti look-ahead strict}

L'anti look-ahead est maintenu à trois niveaux~:

\begin{enumerate}
\item \textbf{Au niveau de chaque sleeve}~: tous les signaux sont calculés avec \code{signal.shift(1)}, de sorte que le signal à l'instant $t$ n'utilise que l'information strictement disponible avant $t$.
\item \textbf{Au niveau du filtre macro}~: l'unemployment rate (publié le premier vendredi du mois suivant) est \emph{forward-filled} pour respecter sa publication réelle. Le yield curve spread est daily et ne nécessite pas de lag.
\item \textbf{Au niveau du vol targeting}~: le levier à $t$ utilise $\max(\mathrm{vol}_{21}, \mathrm{vol}_{63})$ calculé sur les returns $\le t-1$, via un \code{.shift(1)} explicite.
\end{enumerate}

\begin{insightbox}
\textbf{Rappel d'un incident historique}~: la Phase~2A-ERRATA a détecté un biais de look-ahead sur le timeframe 1H qui faisait passer le Sharpe de~$-1.48$ à $+3.10$. Cette expérience a conduit à une politique systématique d'audit \code{.shift(1)} avant toute validation de métrique.
\end{insightbox}

\subsection{Marche à suivre pour reproduire}

\begin{verbatim}
# Prérequis: Python 3.12+, vectorbtpro (licence privée),
#            parquets FX minute et macro dans data/

python -m venv .venv && source .venv/bin/activate
pip install -r requirements.txt

python scripts/generate_phase18_report_artifacts.py   # ~5 min
python scripts/build_phase18_latex_assets.py          # ~2 min
bash scripts/compile_phase18_report.sh                # ~30 s
\end{verbatim}
